\documentclass[11pt]{article}
\usepackage[margin=1in]{geometry}
\usepackage{amsmath,amssymb,amsthm,mathtools}
\usepackage{graphicx}
\usepackage[hidelinks]{hyperref}
\usepackage{microtype}

\newtheorem{theorem}{Theorem}
\newtheorem{lemma}[theorem]{Lemma}
\newtheorem{corollary}[theorem]{Corollary}
\newtheorem{remark}[theorem]{Remark}
\newcommand{\N}{\mathbb N}
\newcommand{\id}{\mathrm{id}}
\newcommand{\dens}{\operatorname{dens}}

\title{A $c_0$-detector criterion for the Banakh--Cauty\\
linear metric AR problem}
\author{Candidate partial-result packet for arXiv:1006.3092}
\date{17 August 2026}

\begin{document}
\maketitle

\noindent\textbf{Status.}
Candidate partial result; likely valid, subject to expert review.  This packet
does not claim to settle either source problem in full.

\section{Source question and result}

Banakh and Cauty ask in Section~4, page~8 of
\cite{BanakhCauty}:
\begin{quote}
\emph{Is each complete linear metric AR-space homeomorphic to a Hilbert
space?}
\end{quote}
They also ask the stronger convex-subset version displayed in the source crop.

\begin{center}
\includegraphics[width=\textwidth]{figures/open_problem_crop.png}
\end{center}

Here a continuous linear map $T:X\to c_0$ is called \emph{compact} if
$T(V)$ is totally bounded for some neighbourhood $V$ of zero.  This is the
definition used by Banakh--Cauty.  The main result of this packet is the
following sufficient condition.

\begin{theorem}[noncompact $c_0$ detector]\label{thm:main}
Let $X$ be a complete linear metric absolute retract.  If there is a
continuous noncompact linear operator
\[
                         T:X\longrightarrow c_0,
\]
then $X$ is homeomorphic to the Hilbert space
$\ell_2(\dens X)$.
\end{theorem}

This removes local convexity from the source theorem whenever such a detector
exists.  For quasi-Banach spaces it gives a particularly transparent partial
answer.

\begin{corollary}[Banach-envelope subcase]\label{cor:envelope}
Let $X$ be a quasi-Banach space which, as a topological space, is an absolute
retract.  If its Banach envelope $\widehat X$ is infinite-dimensional, then
$X\cong\ell_2(\dens X)$.

If instead $\widehat X$ is finite-dimensional and $X$ is infinite-dimensional,
then
\[
                     X\cong N\times \widehat X
\]
linearly and topologically, where $N$ is a complete linear metric AR-space
with trivial continuous dual and trivial Banach envelope.  Thus, within the
quasi-Banach class, the part not reached by the theorem reduces to the
trivial-envelope case.
\end{corollary}

\section{\texorpdfstring{\hspace{0.25em}Proof intuition}{Proof intuition}}

The source proof uses local convexity twice.  The first use is to replace the
identity by a map whose image is locally contained in a finite convex hull.
For an AR this can instead be done topologically: embed $X$ as a closed subset
of a Banach space, take a cover-controlled barycentric approximation there,
and retract it to $X$.  The resulting map is joined to the identity by a small
homotopy and has locally relatively compact image.

The second part of the source argument is linear rather than locally convex.
Noncompactness of $T$ supplies small markers whose $c_0$ images have a large
new coordinate and tiny later-coordinate interactions with all earlier
markers.  A two-marker partition of unity inserts these markers into the
locally precompact approximation.  For two distinct inserted codes, one of
five exhaustive coefficient/order cases exposes a coordinate separated from
zero.  This gives the separated approximation property (SAP), hence LFAP by
Banakh--Cauty, and their Hilbert-space recognition theorem finishes the proof.

The original problem remains open for spaces with no such linear detector.
For a quasi-Banach space, every map to a Banach space factors through its
Banach envelope; consequently the trivial-envelope class is a genuine hard
boundary of this method, not merely a missing estimate.

\section{Proof}

We use two lemmas.  The first is the point at which the AR hypothesis replaces
local convexity.

\begin{lemma}[controlled locally precompact approximation]\label{lem:arapprox}
Let $(X,d)$ be a completely metrizable AR, where $d$ is any compatible
metric.  For every $\varepsilon>0$ there are a map $f:X\to X$ and a homotopy
$H$ from $\id_X$ to $f$ such that every track $H(\{x\}\times[0,1])$ has
$d$-diameter less than $\varepsilon$, and every $x\in X$ has a neighbourhood
$O_x$ for which $f(O_x)$ is relatively compact.
\end{lemma}

\begin{proof}
Choose a bounded complete compatible metric on $X$ and use the Kuratowski
embedding to identify $X$ with a closed subset of a Banach space $E$.  Since
$X$ is an AR, there is a retraction $r:E\to X$.

For each $x\in X$, continuity of $r$ gives a convex norm-ball $W_x$ about
$x$ such that $r(W_x)$ has $d$-diameter less than $\varepsilon$.  Take a
locally finite star-refinement $\mathcal V$ of the cover
$\{W_x\cap X:x\in X\}$, choose $x_V\in V$ for every $V\in\mathcal V$, and
let $(\lambda_V)_{V\in\mathcal V}$ be a subordinate partition of unity.  Put
\[
 g(y)=\sum_{V\in\mathcal V}\lambda_V(y)x_V\quad\hbox{in }E,
 \qquad f=r\circ g,
\]
and
\[
 H(y,t)=r\big((1-t)y+t g(y)\big).
\]
If $\lambda_V(y)>0$, all the participating $V$'s meet at $y$; the
star-refinement therefore puts $y$, all corresponding $x_V$, their convex
hull, and the joining segment inside a single $W_x$.  Hence the homotopy has
the required control.  Local finiteness gives a neighbourhood $O_y$ on which
only finitely many $\lambda_V$ occur.  Thus $g(O_y)$ lies in the convex hull
of finitely many points, whose closure is compact, and $f(O_y)$ is relatively
compact.
\end{proof}

The next lemma isolates and audits the marker part of the argument in
\cite[proof of Lemma~2]{BanakhCauty}.  We include the details because the
published formulas contain harmless index slips and because no local
convexity is used here.

\begin{lemma}[$c_0$ marker lemma]\label{lem:marker}
Let $(X,d)$ be a linear metric space.  Suppose that

\begin{enumerate}
\item for every $\varepsilon>0$, the identity has an $\varepsilon$-small
homotopic approximation $f:X\to X$ which is locally relatively compact; and
\item there is a continuous noncompact linear map $T:X\to c_0$.
\end{enumerate}

Then $(X,d)$ has SAP.
\end{lemma}

\begin{proof}
Fix $\varepsilon>0$.  Choose a balanced zero-neighbourhood $B$ so small that
$B+B$ has $d$-diameter less than $\varepsilon/3$, and, shrinking $B$ if
necessary, assume $T(B)$ is contained in the unit ball of $c_0$.  Since $T$
is noncompact, $T(B)$ is not totally bounded.  The usual compactness
criterion in $c_0$ therefore supplies $0<\delta\le1$ such that
\begin{equation}\label{eq:tailfailure}
 \text{for every }N\text{ there are }u\in B, q\ge N
 \text{ with } |(Tu)_q|\ge\delta.
\end{equation}
Put $\alpha=\delta^3/100$.  Inductively choose $u_r\in B$ and strictly
increasing coordinates $q_r$ such that
\begin{equation}\label{eq:markers}
 |(Tu_r)_{q_r}|\ge\delta,qquad
 |(Tu_s)_{q_r}|<\alpha\quad(s<r).
\end{equation}
Indeed, after $u_0,\ldots,u_{r-1}$ have been chosen, their finitely many
$c_0$ tails are uniformly below $\alpha$, and
\eqref{eq:tailfailure} chooses the next coordinate and marker.

Let $\xi:\N^2\to\N$ be injective, increasing in its first argument, and
satisfying $\xi(i,k)>i$.  This is obtained by dividing $\N$ into countably
many infinite sets.  Write
\[
 u_{i,k}=u_{\xi(i,k)},\qquad
 \phi_{i,k}=e^*_{q_{\xi(i,k)}}T.
\]
Thus $|\phi_{i,k}(u_{i,k})|\ge\delta$; a functional belonging to a later
$\xi$-rank is smaller than $\alpha$ on every earlier marker; and all marker
evaluations have modulus at most one.

Apply Lemma~\ref{lem:arapprox}, with control $\varepsilon/3$, to obtain $f$.
For $s\in\N$, let $C_s$ be the open set of points having a neighbourhood $O$
on which
\[
             |e^*_{q_r}Tf(y)|<\alpha
             \quad(y\in O,\ r\ge s).
\]
Local relative compactness of $f$ and uniform tail decay on compact subsets
of $c_0$ show that the increasing open sets $C_s$ cover $X$.

A standard shrinking of this increasing cover gives increasing open sets
$U_0=\varnothing,U_1,\ldots$ covering $X$ with
$\overline U_{i+1}\subset C_i\cap U_{i+2}$.  Take a partition of unity
$(\lambda_i)$ subordinate to
\[
                  U_{i+1}\setminus\overline U_{i-1}
                  \quad (U_{-1}=U_0=\varnothing).
\]
If $x\in U_{i+1}\setminus U_i$, exactly the two terms $\lambda_i(x)$ and
$\lambda_{i+1}(x)$ can be nonzero, and their sum is one.  Define
\begin{equation}\label{eq:fk}
 f_k(x)=f(x)+\sum_i\lambda_i(x)u_{i,k}.
\end{equation}
The sum is locally finite and in fact has at most two terms.  The straight
homotopy from $f$ to $f_k$, concatenated with the homotopy from $\id_X$ to
$f$, is an $\varepsilon$-homotopy: every difference along its second part is
in $B+B$.

It remains to prove uniform separation.  Fix $k\ne n$ and points
$x\in U_{i+1}\setminus U_i$, $y\in U_{j+1}\setminus U_j$.  Put
\[
 a_0=\xi(i,k),\quad a_1=\xi(i+1,k),\quad
 b_0=\xi(j,n),\quad b_1=\xi(j+1,n),
\]
and write $s=\lambda_{i+1}(x)$, $t=\lambda_{j+1}(y)$.  The four ranks are
distinct, $a_0<a_1$, $b_0<b_1$, and we may interchange the two points so
that $a_1<b_1$.  At any coordinate selected below, its rank is at least
$\max\{i,j\}$; hence the two base terms $Tf(x),Tf(y)$ each contribute less
than $\alpha$.

The following five cases are exhaustive.  In the indicated coordinate, the
displayed lower bound follows from \eqref{eq:markers}; an ``early-marker''
error costs $\alpha$, while a later marker costs at most its coefficient.
\begin{enumerate}
\item If $t>\delta^2/10$, use rank $b_1$ and obtain
$t\delta-5\alpha>\delta^3/20$.
\item If $t\le\delta^2/10$ and $b_0>a_1$, use $b_0$ and obtain
$(1-t)\delta-t-4\alpha>\delta^3/20$.
\item If $t\le\delta^2/10$, $b_0<a_1$, and $s>\delta/4$, use $a_1$ and
obtain $s\delta-t-4\alpha>\delta^3/20$.
\item If $t\le\delta^2/10$, $b_0<a_1$, $s\le\delta/4$, and $a_0<b_0$,
use $b_0$ and obtain $(1-t)\delta-s-t-3\alpha>\delta^3/20$.
\item In the remaining case $a_0>b_0$, use $a_0$ and obtain
$(1-s)\delta-s-t-3\alpha>\delta^3/20$.
\end{enumerate}
For completeness, the inequalities use only $0<\delta\le1$,
$\alpha=\delta^3/100$, $t\le\delta^2/10$, and $s\le\delta/4$ where those
hypotheses occur.  Thus
\[
              \|T(f_k(x)-f_n(y))\|_\infty>\delta^3/20
\]
in every case.  By continuity of $T$, some $\eta>0$, independent of
$x,y,k,n$, satisfies
\[
              d(f_k(x),f_n(y))\ge\eta\qquad(k\ne n).
\]
The images $f_k(X)$ are separated, proving SAP.
\end{proof}

\begin{proof}[Proof of Theorem~\ref{thm:main}]
Lemma~\ref{lem:arapprox} and Lemma~\ref{lem:marker} give SAP.  Banakh and
Cauty prove that every metric space with SAP has LFAP
\cite[Lemma~1]{BanakhCauty}.  Their Banakh--Zarichnyy criterion states that a
convex subset of a linear metric space is homeomorphic to an
infinite-dimensional Hilbert space exactly when it is a completely
metrizable AR with LFAP \cite[Theorem~4]{BanakhCauty}.  Apply it to the
linear space $X$, viewed as a convex subset of itself.  A noncompact operator
cannot have finite-dimensional domain, and density is a topological
invariant, so the Hilbert space is $\ell_2(\dens X)$.
\end{proof}

\begin{proof}[Proof of Corollary~\ref{cor:envelope}]
Let $J:X\to\widehat X$ be the canonical envelope map.  If $\widehat X$ is
infinite-dimensional, the Josefson--Nissenzweig theorem
\cite{Nissenzweig} gives norm-one functionals
$\psi_m\in\widehat X^*$ converging weak-star to zero.  Hence
\[
                    T(x)=(\psi_m(Jx))_{m\in\N}
\]
defines a continuous linear map $X\to c_0$.  It is noncompact.  Indeed, the
envelope unit ball is the closed convex hull of $J(B_X)$, so
$\|\psi_m\|=1$ implies
$\sup_{x\in B_X}|\psi_m(Jx)|=1$.  Therefore $T(B_X)$ has no uniformly
vanishing tails and is not totally bounded.  If the image of any
zero-neighbourhood were totally bounded, scaling a quasi-norm ball into that
neighbourhood would make $T(B_X)$ totally bounded as well.  Thus $T$ is
noncompact, and Theorem~\ref{thm:main} applies.

Now suppose $\widehat X$ is finite-dimensional.  The dense linear subspace
$J(X)$ equals $\widehat X$.  Choose preimages of a basis to obtain a
continuous linear section $S:\widehat X\to X$.  Thus
$X=N\oplus S(\widehat X)$, where $N=\ker J$ is closed and complemented.
It is complete and, as a retract of the AR $X$, is an AR.  Every continuous
functional on $N$ would extend to $X$ through the continuous projection onto
$N$, but all functionals on $X$ vanish on $\ker J$; hence $N^*=\{0\}$ and
its Banach envelope is trivial.
\end{proof}

\section{Upgrade attempts and obstruction audit}

Five materially different upgrade steps were pursued.
\begin{enumerate}
\item The source's locally convex barycentric approximation was replaced by
the AR retraction construction of Lemma~\ref{lem:arapprox}.  This succeeded.
\item The source marker proof was re-derived with only a noncompact
$T:X\to c_0$ and audited by the five exhaustive cases above.  This succeeded.
\item Josefson--Nissenzweig was applied to the Banach envelope.  This upgrades
the theorem to every quasi-Banach AR with infinite-dimensional envelope.
\item A finite-dimensional envelope was split off.  This reduces the
remaining quasi-Banach case to a complemented AR with trivial dual and
trivial envelope.
\item Nonlinear $c_0$ detectors and purely metric separated markers were
tested as routes through the last case.  The separation calculation needs
additivity to isolate a marker from the uncontrolled base term
$f(x)-f(y)$.  Moreover every continuous linear map from a trivial-envelope
quasi-Banach space to $c_0$ is zero by the envelope universal property.
Thus the strongest remaining route would require a genuinely nonlinear
replacement for the marker lemma, not a refinement of its constants.
\end{enumerate}
The fifth point is the reason for stopping short of a full claim: the present
method provably has no detector in the residual class.

\section{Verification, novelty search, and limitations}

\begin{itemize}
\item The AR approximation was checked cover-wise: closed Banach embedding,
global retraction, star-refinement, locally finite partition of unity, and
compactness of a finite convex hull.
\item The marker induction uses exactly the uniform-tail characterization of
totally bounded subsets of $c_0$.  All five coefficient/order cases were
checked independently; they cover $a_1<b_1$ and the symmetric alternative.
\item The only external mathematical inputs are the two results quoted from
the source paper (SAP implies LFAP and the Banakh--Zarichnyy recognition
criterion) and Josefson--Nissenzweig for the corollary.
\item Bounded novelty searches (17 August 2026) used the exact source-problem
phrases, ``noncompact linear operator into $c_0$'' with LFAP/SAP/AR, and
``Banach envelope'' with Hilbert homeomorphism/quasi-Banach.  The arXiv
record, the journal source and its citation neighbourhood were checked.  No
later full solution or this detector/envelope criterion was found.  This is
evidence, not a guarantee of novelty.
\item The result does not settle spaces for which every continuous linear
map to $c_0$ is compact.  In particular, it does not settle the
trivial-Banach-envelope quasi-Banach AR case, and it does not by itself answer
the source's more general convex AR-subset problem.
\end{itemize}

\noindent\textbf{Human-review recommendation.}
Review as a substantial candidate partial theorem.  The highest-value checks
are the cover-controlled retraction construction in Lemma~\ref{lem:arapprox},
the open-cover shrinking that yields the two-adjacent partition, and the
rank/coordinate bookkeeping in the five-case marker table.

\begin{thebibliography}{9}
\bibitem{BanakhCauty}
T. Banakh and R. Cauty,
\emph{Topological classification of closed convex sets in Fr\'echet spaces},
Studia Math. \textbf{205} (2011), 9--22; arXiv:1006.3092.

\bibitem{Nissenzweig}
A. Nissenzweig,
\emph{$w^*$ sequential convergence},
Israel J. Math. \textbf{22} (1975), 266--272.
\end{thebibliography}

\end{document}
